% narrative sentences written after reading the results of the run (values from the database and the drill log)
\newcommand{\LiveText}{The picture was taken at 15:02 on simulated Thursday 17 September, while the system was
still running after the 15 measured days. At that moment 27 of the 30 vehicles were online. 37\% of them were
carrying a passenger and 52\% were waiting for a ride. The fleet had finished 245 trips that day and earned
236,296 LKR in fares. No vehicle had an open idle alert, and 8 vehicles had lost money so far that day (live
fares minus fuel and the part of the daily maintenance cost used up so far).}

\newcommand{\BusinessText}{Over the 15 days the fleet made 150,599 LKR net profit per day on average, and the
vehicles carried a passenger in 39.6\% of their online time. The five vehicles with special profiles
(Section~\ref{sec:sources}) behaved as planned, and the pipeline found each of them without being told:
\begin{itemize}
\item \textbf{V007} made 6,199 LKR on day 1 but lost 17,830 LKR on day 15. Its fuel cost rose from 19.0 to
44.0 LKR per km, and its trips fell from 14 to 4 a day. It was first flagged as becoming unprofitable on day 3,
its first day with a loss. From day 4 to day 10 the report listed it as unprofitable only, because each day
had a service bill or a steady loss. From day 11 the loss grew again every day, so the flag came back.
\item \textbf{V013} lost money only on its four service days, when maintenance cost 21,000 to 26,000 LKR. It
was never flagged as becoming unprofitable, because one large bill is not a falling trend.
\item \textbf{V026} made a profit on every day except its one service day, so a profit rule alone would miss
it. But it claimed 184 km a day while the GPS showed only 138 km, and the report listed it every day for this
gap (28 to 40\%).
\item \textbf{V021} and \textbf{V027} made 2,654 LKR a day, less than half of the 6,267 LKR of the 23
vehicles with no special profile or night shift. They were idle for 514 minutes a day (others: 349) and made
9.7 trips (others: 17.0).
\end{itemize}
The pipeline also found one story that we did not plan. \textbf{V030} works at night, when there are few
rides. It made only 588 LKR a day on average and was flagged three times (days 5, 8 and 14). In total the batch
layer raised 9 ``becoming unprofitable'' flags in 15 days, 6 for V007 and 3 for V030, and none for any other
vehicle.}

\newcommand{\ReconText}{On an average day the speed layer missed about two trips, and its fares were 0.42\%
below the exact batch value (never more than 1.11\%). So the live numbers are good enough to act on during the
day, and the batch numbers are the final answer for money.}

\newcommand{\HealthText}{Every DAG run succeeded. Two runs needed the automatic retry of the sensor
\code{wait\_for\_expense\_file}, both times because we recreated a container (Airflow once and PostgreSQL once)
while the sensor was running. The PostgreSQL restart also stopped the Spark streaming job. Docker restarted
it, and its four queries continued from their checkpoints within about 20 seconds.

We used the recompute path for a real reason. We changed the profit rule so that a service day does not count
as a falling trend, and then triggered the DAG for days 1 to 13 through the Airflow REST API with the
\code{sim\_date} parameter. The 13 runs took 7 minutes and 1 second in total, and the ``becoming unprofitable''
flags for those days fell from 20 to 6. A rule change in a Lambda system is only a rerun of the batch layer.

The alerts also caught one real mistake. At the start of the long run we lowered the memory of the Spark
streaming driver, and Spark refused to start because it needs at least 450\,MB. SpeedLayerStale and
ServiceDown (Spark) fired within two minutes of the first error, and both were resolved less than four
minutes after it, once we set the memory back to 512\,MB.}

\newcommand{\DrillText}{Figure~\ref{fig:health} shows the result. Within 12 seconds Prometheus marked
ServiceDown (stream-producer) and NoTelemetryIngested as pending. Both were firing at 73 seconds, because the
rules wait one minute to avoid false alarms. The API webhook received ServiceDown at 81 seconds. It did not
receive NoTelemetryIngested, because the inhibition rule hides it while the producer is known to be down. So
the team gets one clear message instead of two. SpeedLayerStale was firing at 113 seconds and reached the
webhook at 122 seconds. This is correct, because the live dashboards had stopped changing. We started the
producer again at 215 seconds. All alerts cleared within 18 seconds, and the ``resolved'' messages arrived 27
and 47 seconds after the restart. The simulator does not send the events of the time when it was stopped, so
the drill day (16 September) has a gap of about seven simulated hours. For this reason the results in
Section~\ref{sec:daily} stop at 15 September.}
